\documentclass[12pt]{article}
\usepackage[utf8]{inputenc}
\usepackage[T1]{fontenc}
\usepackage{hyperref}
\title{Fore Language Resources for K-12 Curriculum Development}
\author{Compiled by Claude Sonnet 5.5 (claude-sonnet-5-5) for Andr\'as Kornai}
\date{}
\begin{document}
\maketitle

\section{Introduction}
Fore (ISO 639-3 \texttt{for}) is a Trans-New Guinea language spoken in the Okapa District, Eastern Highlands Province, Papua New Guinea. It has three dialects. Only academic description was found in searches.

\section{Primary Educational Resources}
No primers, textbooks or school materials were found.

\section{Linguistic and Reference Materials}
\begin{itemize}
\item Graham Scott, \emph{The Fore Language of Papua New Guinea} (Pacific Linguistics, 1978; from his 1977 ANU thesis), a tagmemic grammar. Open repository record: \url{https://openresearch-repository.anu.edu.au/items/45a1eeba-eb8c-4651-9554-c3957a4655ea}.
\item Bibliographic record in Glottolog: \url{https://glottolog.org/resource/reference/id/107238}.
\item Numeral system data for Fore: \url{https://lingweb.eva.mpg.de/channumerals/Fore.htm}.
\end{itemize}

\section{Religious and Community Resources}
None found. An encyclopedia entry says missionary linguists developed an orthography, but gives no details: \url{https://www.everyculture.com/Oceania/Fore-i-Orientation-i.html}.

\section{Digital and Audio Resources}
None found. No dictionary was located in these searches.

\section{Curriculum Development Context}
The Scott grammar is the main basis for any curriculum. Orthography details, existing literacy materials and a dictionary would need to be sought from SIL Papua New Guinea and local community sources; none were verified here.
\end{document}